\begin{tablalafrox}{Diccionario del Producto 5: Integraciones y datos migrados. Fuente: según cita por cuenta en cada encabezado.}{tab:T14-dic-prod-5}
  {L{1.2cm} Y Y Y C{0.9cm} L{2.3cm}}
  {\cab{Código} & \cab{Paquete de trabajo} & \cab{Entregable} & \cab{Criterio de aceptación} & \cab{Resp.} & \cab{Período}}
  \multicolumn{6}{>{\columncolor{lafroxBg}\RaggedRight\arraybackslash}p{\dimexpr\textwidth-2\tabcolsep\relax}}{\textbf{Cuenta 5.1 — Integraciones con el ERP, el WMS y la flota.} Esta cuenta reúne las integraciones con el ERP, el WMS y la flota, como servicios compartidos por los módulos. {\color{lafroxGris}Fuente: (Bases Administrativas, arts. 17.2–17.6 y Formulario T-14, p. 63); (Bases Técnicas Transversales, secciones 17.2–17.6, pp. 32–34).}} \tabularnewline*
  5.1.1 & Integración con el ERP como único emisor de la guía de despacho (INT-06 e INT-07) & Integración en QA. & Toda guía la emite el ERP; cada entrega se concilia con su guía el mismo día; una caída del ERP no detiene el despacho. & ARQ & E1 · Meses 5 y 6; termina antes de la prueba de integración del H4 (T-15, sección 6.1, programación de paquetes con entregable). \\
  5.1.2 & Convivencia con el WMS de 2013 durante el reemplazo & Mecanismo de convivencia y procedimiento para apagar el WMS. & No hay doble escritura de stock durante la convivencia. & ARQ & E1 · Antes del H4. \\
  5.1.3 & Integración con la telemetría de los 42 camiones propios (INT-15) & Integración en QA. & M12 recibe la posición de los 42 camiones. & ARQ & E1 · Antes del H4. \\
  \midrule
  \multicolumn{6}{>{\columncolor{lafroxBg}\RaggedRight\arraybackslash}p{\dimexpr\textwidth-2\tabcolsep\relax}}{\textbf{Cuenta 5.2 — Integraciones con servicios externos.} Esta cuenta reúne las conexiones con servicios externos identificadas con los códigos de interfaz de la arquitectura del Capítulo 4 (Subdocumento 4). Cada integración define su conducta ante fallas del tercero para mantener la operación. {\color{lafroxGris}Fuente: (Bases Administrativas, arts. 17.2–17.6 y Formulario T-14, p. 63); (Bases Técnicas Transversales, secciones 17.2–17.6, pp. 32–34).}} \tabularnewline*
  5.2.1 & Intercambio de lote y trazabilidad con los proveedores & Intercambio operativo con los proveedores priorizados. & El proveedor de lácteos acepta la evidencia de trazabilidad: un simulacro de retiro de uno de sus lotes en QA, con datos de prueba de recepción, lote y despacho, identifica clientes, fechas, cantidades y guías en menos de 2 horas (R18-01). & DAT & E1 · Antes del H5. \\
  5.2.2 & Integración con la pasarela de pago para el cobro con tarjeta (INT-09) & Integración en QA. & Un cobro de prueba se autoriza, se concilia con su entrega, y una falla de la pasarela no lo deja como aprobado. & DES & E1 · Antes del H4. \\
  5.2.3 & Integración con el servicio de mapas y geocodificación (INT-10) & Integración en QA. & Las rutas se generan con las direcciones georreferenciadas, y la ruta precargada funciona sin señal. & DES & E1 · Antes del H4. \\
  5.2.4 & Integración con el servicio de avisos al cliente (INT-11) & Integración en QA. & Un aviso de prueba llega por el canal configurado, y los reintentos quedan registrados. & DES & E1 · Antes del H4. \\
  \midrule
  \multicolumn{6}{>{\columncolor{lafroxBg}\RaggedRight\arraybackslash}p{\dimexpr\textwidth-2\tabcolsep\relax}}{\textbf{Cuenta 5.3 — Intercambio electrónico con las cadenas.} Esta cuenta reúne el intercambio electrónico con las cadenas, identificado con los códigos de interfaz de la arquitectura del Capítulo 4 (Subdocumento 4). Las conexiones con las cadenas tienen una conducta definida frente a fallas del tercero para mantener la operación. {\color{lafroxGris}Fuente: (Bases Administrativas, arts. 17.2–17.6 y Formulario T-14, p. 63); (Bases Técnicas Transversales, secciones 17.2–17.6, pp. 32–34).}} \tabularnewline*
  5.3.1 & Intercambio electrónico con la principal cadena, con sus pruebas aprobadas (INT-08) & Intercambio certificado por la cadena. & La cadena aprueba por escrito las pruebas antes del H10. El perfil de mensajes, el acceso de prueba y el calendario de certificación se acuerdan desde el levantamiento de interfaces y se registran antes de construir el intercambio. & DES & E2 · Del 1 de octubre de 2027 al 31 de mayo de 2028 (meses 9 a 16). Prueba de punta a punta con 4.4.1; conformidad antes de H10. \\
  5.3.2 & Incorporación de las demás cadenas al intercambio electrónico (INT-08) & Cadenas incorporadas. & Cada cadena tiene conformidad documentada antes de activarse. Si una cadena no obtiene la certificación a tiempo, el Jefe de Proyecto escala la desviación y mantiene la carga manual controlada. & DES & E2 · Del 2 de mayo al 14 de julio de 2028. Las cadenas se certifican antes de cada activación y todo el alcance queda habilitado antes de las cuatro semanas finales de evidencia de la marcha blanca. \\
  \midrule
  \multicolumn{6}{>{\columncolor{lafroxBg}\RaggedRight\arraybackslash}p{\dimexpr\textwidth-2\tabcolsep\relax}}{\textbf{Cuenta 5.4 — Migración de datos.} Esta cuenta reúne la migración de los datos actuales a la solución. Comprende dos ensayos previos a la migración definitiva y una conciliación sin diferencias no explicadas (BTT, núm. 20.1, p. 35). {\color{lafroxGris}Fuente: (Bases Administrativas, arts. 17.2–17.6 y Formulario T-14, p. 63); (Bases Técnicas Transversales, secciones 17.2–17.6, pp. 32–34).}} \tabularnewline*
  5.4.1 & Informe de perfilamiento y datos saneados & Informe de calidad de datos y datos corregidos. & Cada defecto tiene su corrección o una aceptación del dueño del dato. & DAT & E1 · Antes de 5.4.2. \\
  5.4.2 & Datos maestros cargados: 8.400 SKU, clientes e instalaciones & Maestro cargado en Preproducción. & El maestro cuadra con el origen en los dos ensayos. & DAT & E1 · Antes de la primera ola. \\
  5.4.3 & Saldos migrados desde el WMS de Talca & Saldos de Talca migrados y conciliados. & Los documentos y los lotes coinciden, y cada diferencia de inventario queda clasificada y explicada. & DAT & E1 · Antes de la ola de Talca. \\
  5.4.4 & Saldos cargados en Concepción desde el conteo físico & Saldos de Concepción cargados. & Lo cargado coincide con el conteo, y cada diferencia frente a las planillas queda clasificada. & DAT & E1 · Antes de la ola de Concepción. \\
  5.4.5 & Acta de conciliación y corte fuera de las ventanas de congelamiento & Acta de corte por sitio. & No quedan diferencias sin explicar, y la fecha no es una fecha prohibida. & DAT & E1 · Antes del H6. \\
\end{tablalafrox}
